\section{Compact permutations for braid component checks}
\label{sec:compact-permutations}

The lazy factor preparation of Section~\ref{sec:lazy-forest} exposes a
remaining cost: computing a strand permutation for every supplied grammar
rule. A local three-strand subword in a large forest previously occupied a
full $m$-entry tuple. This cost was paid in production and again in independent
replay, even when one finite factor settled the knot verdict. The maintained
implementation now stores local permutations compactly and evaluates only
root-reachable permutation nodes. It still validates every supplied rule.

\subsection{Exact composition with adaptive storage}

Use zero-based strand indices. A generator of either sign exchanges two
adjacent indices. If a concatenation has child permutations $p,q$, the
established convention is
\[
 (p\circ q)(i)=p(q(i)).
\]
A sparse representation stores exactly the pairs $i\mapsto p(i)$ for which
$p(i)\ne i$. An absent key means a fixed point, and the empty map represents
the identity. To compose two sparse operands, start with a copy of $p$.
For every stored pair $i\mapsto j$ of $q$, set its image to $p(j)$, using
$j$ when that key is absent in $p$. Delete the result entry if its image is
$i$. For indices outside the support of $q$, $q(i)=i$, so the initial copy
already has the correct image. This proves the formula for every strand,
including overlapping supports and products that cancel to the identity.

When the resulting map moves more than half the strands, it is converted to
a dense tuple. Dense--dense composition indexes the first tuple by the
second. Dense--sparse composition copies the dense tuple and updates just
the sparse operand's domain. Sparse--dense composition performs one sparse
lookup for each image in the dense tuple. These are the same formula in
three storage cases. A dense descendant remains dense even if its permutation
later becomes sparse or the identity; there is no claim that storage always
tracks minimal support.

All operands are treated as immutable. An identity product can therefore
reuse its nonidentity operand without copying, including when a grammar node
has several parents. Strands counts at most eight retain the simple dense
path to avoid representation dispatch overhead. The half-density and
small-strand thresholds are implementation heuristics, not topological facts
or claims of an optimal policy. The internal dense ablation retains the new
reachability check while disabling compact storage.

\subsection{Reachability, source validation and proof compatibility}

The existing reverse root-weight pass already assigns zero to unreachable
nodes. A live concatenation has two live children because both child weights
receive a positive contribution, including twice the contribution if the
children coincide. Thus a forward pass can skip every zero-weight permutation
without losing any dependency of the root. Unreachable rules still undergo
full syntax, index, generator and metadata validation. They remain part of
the exact source grammar and serialized-byte accounting; this is not source
pruning or a change in certificate identity.

Before allocating any strand-sized representation, the validator checks for
a missing generator label in the root's unsigned occurrence counts. A missing
label separates the strand intervals and proves that the closure is a link.
This guard works with a binary strand count far larger than the supplied
rule count. If there is no gap, every one of the $m-1$ labels occurs in some
terminal, so $m\le s+1$ for $s$ supplied rules. The strand-sized allocations
are made only after this bound is established.

Finally the root is materialized as the same full permutation list as before.
Its orbit from strand zero must contain all $m$ strands to establish a knot.
This check remains before factor scheduling or any selected negative-factor
proof. A link containing a nontrivial component cannot acquire a nontrivial
\emph{knot} verdict from that component. Producer and replay use the same
validated permutation arithmetic but replay still reconstructs it afresh
from the exact source, without trusting a claimed component count.

No public option or certificate schema changes. Exact summary fields,
canonical proofs, factor ordering, verdicts and charged abstract work remain
unchanged on the differential audit. The callback still runs at each supplied
rule in the permutation pass, even when its permutation is skipped. The work
counter counts cooperative operations, not every tuple entry or hash probe;
a faster or smaller representation need not reduce that counter. Work and
time limits still yield no partial certificate, and external cancellation
propagates.

\subsection{Cost and remaining asymptotic limitations}

Let $s$ be the supplied rule count and $m$ the number of strands. For a live
node $v$, write $r_v$ for its representation size: the moved-point count for
a sparse map, and $m$ for a dense tuple. A sparse--sparse composition performs
$O(r_u+r_w)$ dictionary operations before any dense conversion. A conversion
costs $O(m)$ operations but occurs only after more than $m/2$ output entries
have been constructed. A composition with a dense operand costs $O(m)$,
which is also $O(r_u+r_w)$. Live terminals use two entries in the compact path;
identity aliases allocate no permutation entries.

Thus, with constant-cost dictionary operations, the permutation pass uses
\[
 O\!\left(s+m+\sum_{v=(u,w)\text{ live}}(r_u+r_w)\right)
\]
word operations and at most $O(s+m+\sum_{v\text{ live}}r_v)$ stored entries.
The sums bound storage even when aliases share objects. The ordinary dense
bound $O(sm)$ remains the worst case in this model, while local forest nodes
can have representation size independent of the total strand count.
The current implementation retains node representations until the root has
been evaluated; it does not claim optimal lifetime management.

These are representation and dictionary-operation bounds, not an unconditional
bit-cost theorem for Python hash tables. Python's conventional amortized
hash-table cost model gives the stated elapsed-work interpretation. Charging
up to $O(m)$ probes per lookup instead gives a conservative polynomial
$O(sm^2)$ probe bound. Strand keys and images have $O(\log m)$ bits; the
separate validated lengths, exponents and root weights must still be charged
at their own binary lengths. No deterministic subexponential or
quasi-polynomial theorem is inferred from a hash-table timing result.

In particular, a large dense descendant need not become compact again after
cancellation. All-factor positive recognition still pays every required
projection and leaf proof. The input remains a supplied braid grammar; this
change neither constructs one of controlled size from an arbitrary knot
diagram nor bounds the exceptional-factor parameter on every diagram. The
restricted recognition theorems of the preceding sections are unchanged.

\subsection{Native differential validation and measurements}

Eight focused tests pass in 0.196 seconds. They exhaust all 576 ordered pairs
of four-strand permutations in each of the four sparse/dense representation
combinations, checking results and operand immutability. A separate oracle
expands 400 bounded random shared grammars and applies adjacent swaps directly.
Other cases check repeated child references, a $2^{256}$-power grammar,
cancellation, dead-rule validation, a strand count of $10^{100}$, complete
component checks on links, exact public work ceilings and proof equality with
the dense ablation. The small-strand path is tested with compact composition
disabled.

The native compatibility audit loads the actual previous compressed-braid
package at commit \texttt{d0f0e37761b5}. It combines 240 prior cases, twelve
benchmark sources and 150 additional seeded sources with shared reachable
and unreachable powers. With the exponential fallback disabled equally in
both versions, the 402 cases give 103 unknot, 132 nontrivial-knot, 148 link
and nineteen inconclusive results. All public result fields and root summaries
agree exactly, and every old and new certificate replays in both versions.
The 4.117-second audit retains its complete sources and hashes in
\path{data/compact-permutations-audit.json}.

The full maintained suite passes 948 tests with Regina in 222.385 seconds.

The benchmark has twelve fixtures, five arms, one shuffled excluded warmup
and seven shuffled measured rounds per fixture: 420 measured calls and sixty
warmups. The arms are the pinned old package, an identical old control,
current defaults, an identical current control, and current code with dense
permutation storage forced internally. Policy binding for this last arm is
outside the timer; all source validation, permutation construction, factor
recognition, proof transport accounting and mandatory replay are inside.
Every arm produces the same complete public result. Separate untimed
\texttt{tracemalloc} calls measure peak traced allocations of a complete root
summary, including its validation and metadata. These are not whole-process
resident-memory measurements or enforced memory ceilings.

\begin{center}\small
\begin{tabular}{@{}lrrrrr@{}}
\toprule
Input & old ms & new ms & dense ms & ratio & peak KiB old/new \\
\midrule
late\_finite\_8 & 11.967 & 11.155 & 12.085 & 1.072 & 204.0/184.9 \\
late\_finite\_24 & 43.256 & 33.085 & 42.199 & 1.310 & 1292.0/578.3 \\
late\_finite\_64 & 188.472 & 89.181 & 178.327 & 2.103 & 8002.7/1674.1 \\
early\_compressed & 109.791 & 99.178 & 108.404 & 1.115 & 1257.1/574.8 \\
positive\_forest & 118.311 & 115.619 & 117.032 & 1.023 & 609.1/374.2 \\
all\_singleton & 13.562 & 12.305 & 13.787 & 1.103 & 288.4/187.6 \\
single\_wider & 1.212 & 1.235 & 1.242 & 0.977 & 2.5/2.6 \\
direct\_three & 4.102 & 4.125 & 4.120 & 0.994 & 11.6/11.7 \\
dense\_random\_link & 49.501 & 34.507 & 48.795 & 1.432 & 2492.5/1578.6 \\
dense\_power\_link & 19.512 & 16.999 & 18.060 & 1.139 & 968.4/874.8 \\
unreachable\_rules & 43.521 & 32.854 & 33.371 & 1.302 & 1236.1/137.1 \\
huge\_strand\_gap & 0.088 & 0.088 & 0.090 & 1.003 & 0.5/0.5 \\
\bottomrule
\end{tabular}
\end{center}


The ratio is the median paired old/new time. The dense column still skips
unreachable permutation nodes, so it separates compact storage from that
optimization. Raw source grammars, sample orders and times, resource counts,
proof sizes and source hashes are retained in
\path{../fast/results/compact_permutations_20261008.json}. The complete run,
including separate memory measurements, takes 27.469 seconds. Identical-old
paired ratios range from 0.972 to 1.017; identical-current ratios range from
0.989 to 1.021. Small changes near this range are not evidence of a gain.

On the 64-positive-factor forest followed by a finite negative factor,
complete recognition falls from 188.472 to 89.181 ms, with paired ratio
2.103. The dense ablation takes 178.327 ms, locating most of the gain in
compact representations. Root-summary peak traced allocation falls from
8,194,752 to 1,714,252 bytes, a factor of 4.780. Both public calls still charge
98,011 abstract work steps and produce the same proof. The eight- and
24-positive-factor cases have paired ratios 1.072 and 1.310 respectively;
the smaller case also has a smaller memory reduction.

The positive sixteen-factor forest needs all leaf work and has paired ratio
1.023, close to the control range; no systematic whole-call improvement is
claimed. Its root-summary allocation nevertheless falls from 623,728 to
383,220 bytes. The early compressed-negative case has paired ratio 1.115;
the singleton chain has ratio 1.103. The small wider-factor and direct
three-braid controls have ratios 0.977 and 0.994, so the compact policy is
not presented as universally faster.

The dense random-link control improves from 49.501 to 34.507 ms, with ratio
1.432. Although long prefixes become dense, their terminal operands remain
sparse, permitting copy-and-update composition. The repeated dense-power
control has ratio 1.139 but only a 1.107-fold memory reduction, illustrating
that dense descendants remain dense after cancellation. In the unreachable
rule control, old/new medians are 43.521/32.854 ms and the dense ablation is
33.371 ms: here skipping dead permutations explains most of the gain. Its
root-summary peak allocation falls from 1,265,744 to 140,368 bytes. The enormous
strand-count gap control remains constant-size and has no timing gain.

\paragraph{Scope of the advance.}
Local permutation support now reduces a measured preparation bottleneck in
both proof production and replay. The finite obstruction, three-braid and
exceptional-cube mathematics is unchanged. Establishing a complete source-bound
conversion and controlling its size for arbitrary knot diagrams remains
necessary before these restricted algorithms can imply a general
quasi-polynomial recognition theorem.
